\documentclass[10pt,a4paper]{amsart}
\usepackage[margin=19mm]{geometry}
\usepackage{amsmath,amssymb,mathtools}
\usepackage[scr=rsfs,cal=euler]{mathalfa}
\usepackage{xfrac}
\usepackage[unicode,hidelinks,bookmarks=false]{hyperref}
\newcommand{\ie}{i.e.\ }
\newcommand{\cf}{cf.\ }
\newcommand{\resp}{respectively\ }
\newcommand{\Subsection}{\subsection}
\providecommand{\nobreakdash}{\nobreak\hbox{-}}
\setlength{\emergencystretch}{2em}
\multlinegap0pt
\usepackage{mathtools}
\usepackage{xfrac}
\usepackage{stackengine}
\usepackage{tikz}
\usepackage{booktabs}
\usepackage{multirow}
\usepackage{enumerate}
\usepackage{mdframed}
\usepackage{nicefrac}
\usepackage[normalem]{ulem}
\usepackage{cancel}
\usepackage[all]{xy}
\usepackage{algorithm}\usepackage{algpseudocode}
\usepackage{bbm}
\usepackage[shortlabels]{enumitem}
\newtheorem{theorem}{Theorem}
\def\RR{\mathbb{R}}
\def\ba{\mathbf{a}}
\def\be{\mathbf{e}}
\def\bx{\mathbf{x}}
\def\bzero{\mathbf{0}}
\title{l'H\^opital rules for complex-valued functions in higher dimensions}
\author{Albert Chern, Sadashige Ishida}
\date{Source: arXiv:2602.09958v1 (CC BY 4.0)}
\begin{document}
\maketitle

\noindent\emph{Source and licence.} l'H\^opital rules for complex-valued functions in higher dimensions. Version arXiv:2602.09958v1 (CC BY 4.0), distributed by arXiv under the Creative Commons Attribution 4.0 International licence, http://creativecommons.org/licenses/by/4.0/, as recorded in the arXiv licence metadata for this exact version and shown on the abstract page https://arxiv.org/abs/2602.09958v1. Full licence notice: \texttt{licenses/arxiv-2602.09958-CC-BY-4.0.txt}.\par\smallskip

\noindent\textit{Editorial scope.} Counted unit: the article's theorem thm:RealQuotientnD (source0.tex lines 796--799). The article's complete original proof of it is delivered with it (source0.tex lines 1152--1163, one \texttt{proof} environment of 12 lines, which the article places away from the statement). No mathematics is written, repaired or re-derived: the statement and the proof are the article's own lines, in the article's own notation, and no retained line is substituted or re-flowed.  Retained uncounted as the source-local context the proof needs: lines 1098--1101.  Nothing was omitted from any retained span, so the package carries no omission record.  No retained line contains a citation, so the wrapper carries no bibliography and no bibliography entry.  No retained line contains a figure environment, an image-inclusion command or drawing code, so no image is delivered and the package declares no asset dependency.  Editorial formatting only, in addition: an AMS \texttt{amsart} preamble that restates the article's own declarations and macros used by the retained lines, namely the environments \texttt{theorem} on separate counters, together with the standard packages those lines need.  The retained environments print in this document as Theorem 1 in order of appearance, whereas the article prints the same items with the numbering of its own full document.  The complete native source of the article as distributed in the arXiv e-print for this exact version is bundled unchanged as \texttt{sources/194360/source0.tex}.
\begin{align}
\label{eq:WhitneyTowardsMultidiemsion}
        f(x_1,x_2,\ldots,x_n) = f(0,x_2,\ldots,x_n) + x_1 g_1(x_1,x_2,\ldots,x_n)
\end{align}

\begin{theorem}[Smooth quotient of real functions in arbitrary dimension]
\label{thm:RealQuotientnD}
    Let \(f,g\in C^\infty(\Omega)\) be smooth real-valued functions with simple zeros and a common zero set \(\Sigma_{f} = \Sigma_{g}\).  Then there exists a non-vanishing smooth function \(\varphi\in C^\infty(\Omega)\) such that \(f = g\varphi\).
\end{theorem}

\begin{proof}[Proof of \autoref{thm:RealQuotientnD}]
    Let \(\Sigma\) be the common zero set of \(f,g\in C^k(\Omega)\).  On \(\Omega\setminus\Sigma\) define \(\bx\in\Omega\setminus\Sigma\) define \(\varphi(\bx) = f(\bx)/g(\bx)\), which belongs to \(C^{k}(\Omega\setminus\Sigma)\) and is non-vanishing.
    It remains to show that \(\varphi\) is \(C^{k-1}\)-continuously defined in a neighborhood of any point of \(\Sigma\).  Let \(\ba\in\Sigma\).  By the implicit function theorem, there exists a neighborhood \(U\) of \(\ba\), a neighborhood \(\tilde U\subset\RR^n\) of the origin \(\bzero\in\RR^n\), and a \(C^k\) bijection \(\Phi\colon \tilde U\xrightarrow{\simeq} U\) such that \(\Phi(\bzero) = \ba\) and \(\tilde\Sigma = \Phi^{-1}(\Sigma)\) is the coordinate hyperplane \(\tilde\Sigma = \{(x_1,\ldots, x_n)\,|\,x_1 = 0\}\cap\tilde U\).
    Define \(\tilde f = f\circ\Phi\) and \(\tilde g = g\circ\Phi\), both of which are \(C^k\) functions on \(\tilde U\) and have simple zeros on \(\tilde\Sigma\).  
    In this coordinate system, \(\nabla \tilde f(0,x_2,\ldots,x_n) = {\partial \tilde f\over\partial x_1}(0,x_2,\ldots,x_n)\be_1\) and \(\nabla \tilde g(0,x_2,\ldots,x_n) = {\partial \tilde g\over\partial x_1}(0,x_2,\ldots,x_n)\be_1\), where \(\be_1\) is the coordinate vector \((1,0,\ldots,0)\), and \({\partial\tilde f\over\partial x_1}\) and \({\partial\tilde g\over\partial x_1}\) are nonvanishing \(C^{k-1}\) function on \(\tilde\Sigma\).
    By \eqref{eq:WhitneyTowardsMultidiemsion}, there exists \(p,q\in C^{k-1}(\tilde U)\) such that for \((x_1,\ldots,x_n)\in \tilde U\)
    \begin{align*}
        \tilde f(x_1,\ldots,x_n) = x_1 p(x_1,\ldots,x_n),\quad\tilde g(x_1,\ldots,x_n) = x_1q(x_1,\ldots,x_n)
    \end{align*}
    and \(p = {\partial\tilde f\over\partial x_1}\neq 0\) and \(q = {\partial\tilde g\over\partial x_1}\neq 0\) on \(\tilde\Sigma\).
    Therefore there exists a neighborhood of \(\tilde\Sigma\), without loss of generality \(\tilde U\), on which \(p\) and \(q\) are nonvanishing.  Therefore, \({\tilde f\over\tilde g} = {p\over q}\) is a nonvanishing \(C^{k-1}\) continuous function on \(\tilde U\) and agrees with \(\varphi\circ \Phi\) on \(\tilde U\).  Thus \(\varphi\) extends to a \(C^{k-1}\)-continuous function in a neighborhood of \(\ba\).
\end{proof}

\end{document}
